\chapter{Analytical geometry}
Analytical geometry, also called coordinate or Cartesian geometry, is the study of properties, relationships and measurement of points, lines and angles in the Cartesian plane. Geometrical shapes are defined using a coordinate system and algebraic principles. Some consider the introduction of analytical geometry to be the beginning of modern mathematics.\par 

\section{Drawing figures on the Cartesian plane}
If we are given the coordinates of the vertices of a figure, we can draw the figure on the Cartesian plane. For example, quadrilateral $ABCD$ with coordinates $A(1;1)$, $B(3;1)$, $C(3;3)$ and $D(1;3)$  

\setcounter{subfigure}{0}
\begin{figure}[H] % horizontal\label{m39107*id63458}
\begin{center}
\scalebox{.8}{
\begin{pspicture}(-5,-5)(5.5,5.5)
% \psaxes{<->}(0,0)(5,5)
% \psgrid[gridcolor=lightgray,linecolor=lightgray,subgriddiv=1](0,0)(0,0)(4,4)
\psaxes[linewidth=2pt,labels=none,ticks=none]{->}(0,0)(0,0)(5,5)
\pspolygon[linewidth=.1cm](1,1)(1,3)(3,3)(3,1)(1,1)
\uput[dl](1,1){\Large{$A$}}
\uput[dr](3,1){\Large{$B$}}
\uput[ur](3,3){\Large{$C$}}
\uput[ul](1,3){\Large{$D$}}
\uput[l](5.7,0){\Large{$x$}}
\uput[d](0,5.7){\Large{$y$}}
\end{pspicture}
}
% \caption{Drawing a figure on the Cartesian plane}
\end{center}
\label{fig:cartesianplane}
\end{figure} 

The order of the letters for naming a figure is important. It indicates the order in which points must be joined: $A$ to $B$, $B$ to $C$, $C$ to $D$ and $D$ back to $A$. It would also be correct to write quadrilateral $CBAD$ or $BADC$ but it is better to follow the convention of writing letters in alphabetical order.     

\section{ Distance between two points}
\Definition{Point}{A point is an ordered pair of numbers written as $(x;y)$.}
A point is a simple geometric object having location as its only property. 
\Definition{Distance}{Distance os a measure of the length between two points.}

\begin{activity}{Distance between two points}
Points $P~(2;1)$, $Q~(-2;2)$ and $R~(2,-2)$ are given. 
\begin{itemize}
 \item Can we assume that $\hat{R}=90^{\circ}$? If so, why?
\item Apply the Theorem of Pythagoras in $\triangle PQR$ to find the length of $PQ$.
\end{itemize}


\setcounter{subfigure}{0}
\begin{figure}[H] % horizontal\label{m39107*id63458}
\begin{center}
\scalebox{.8}{
\begin{pspicture}(-5,-5)(5.5,5.5)
% \psaxes{<->}(0,0)(5,5)
% \psgrid[gridcolor=lightgray,linecolor=lightgray,subgriddiv=1](0,0)(-3,-3)(3,3)
\psaxes[linewidth=2pt,labels=none,ticks=none]{->}(0,0)(-3,-3)(3,3)
\psline[linewidth=0.04cm](-2,-2)(2,-2)(2,1)
\psline[linestyle=dashed,linewidth=0.04cm](-2,-2)(2,1)
\uput[l](-2,-1.8){\Large{$Q$}}
\uput[dl](-1.5,-2){\Large{$(-2;-2)$}}
\uput[dr](2,-2){\Large{$R$}}
\uput[u](2.3,0.5){\Large{$P$}}
\uput[ur](2,1){\Large{$(2;1)$}}
\uput[l](4,0){\Large{$x$}}
\uput[d](0,4){\Large{$y$}}
\uput[dr](2.5,-2){\Large{$(2;-2)$}}
\end{pspicture}
}
% \caption{Triangle $PQR$}
\end{center}
\label{fig:trianglePQR}
\end{figure} 
\end{activity}       
% SOLUTION FOR TEACHERS' GUIDE BELOW:
% In figure~\ref{fig:trianglePQR}, it can be seen that the length of the line $PR$ is 3 units and the length of the line $QR$ is four units. However, $\Delta PQR$, has a right angle at $R$. Therefore, the length of the side $PQ$ can be obtained by using the Theorem of Pythagoras:\par 
%       
% \begin{eqnarray*}
% P{Q}^{2} & = & P{R}^{2}+Q{R}^{2} \\ \
% \therefore P{Q}^{2} & = & {3}^{2}+{4}^{2} \\ 
% \therefore PQ & = & \sqrt{{3}^{2}+{4}^{2}}=5  
% \end{eqnarray*}
% The length of $PQ$ is the distance between the points $P$ and $Q$.\par 
To derive a general formula for the distance between two points $A~({x}_{1};{y}_{1})$ and $B~({x}_{2};{y}_{2})$, we use the Theorem of Pythagoras:\par 

\setcounter{subfigure}{0}
\begin{figure}[H] % horizontal\label{m39107*id63458}
\begin{center}
\scalebox{.8}{
\begin{pspicture}(-5,-5)(5.5,5.5)
% \psaxes{<->}(0,0)(5,5)
% \psgrid[gridcolor=lightgray,linecolor=lightgray,subgriddiv=1,gridlabels=0.0cm](0,0)(-3,-3)(3,3)
\psaxes[linewidth=1pt,labels=none,ticks=none]{<->}(0,0)(-3,-3)(3,3)
\psline[linewidth=0.04cm](-2,-2)(2,-2)(2,1)
\psline[linestyle=dashed,linewidth=0.04cm](-2,-2)(2,1)
\uput[l](-2,-1.8){\Large{$A$}}
\uput[dl](-1,-2){\Large{$(x_{1};y_{1})$}}
\uput[dr](2,-1.8){\Large{$C$}}
\uput[ur](2.5,-2.5){\Large{$(x_{2};y_{1})$}}
\uput[u](2.3,0.5){\Large{$B$}}
\uput[ur](1.8,1){\Large{$(x_{2};y_{2})$}}
\uput[l](4,0){\Large{$x$}}
\uput[d](0,4){\Large{$y$}}
\uput[d](-0.2,-0.05){\Large{$0$}}
\end{pspicture}
}
\end{center}
\end{figure}       
}
\Tip{Always draw a sketch - it helps with your calculation and makes it easier to check if your answer is correct.}
\begin{align*}
AB^2&=AC^{2}+BC^{2}\\
\therefore AB&=\sqrt{AC^{2}+BC^{2}}
\end{align*}
}And


\begin{eqnarray*}
AC & = & {x}_{2}-{x}_{1}\\
BC & = & {y}_{2}-{y}_{1} 
\end{eqnarray*}
Therefore,
\begin{eqnarray*} 
AB & = & \sqrt{A{C}^{2}+B{C}^{2}} \\ 
& =& \sqrt{(x_{2}-x_{1})^{2}+(y_{2}-y_{1})^{2}} 
\end{eqnarray*}
Therefore to calculate the distance between any two points, $({x}_{1};{y}_{1})$ and $({x}_{2};{y}_{2})$, we use:\par 
\Identity{Distance}{$\mbox{Distance}=\sqrt{{({x}_{1}-{x}_{2})}^{2}+{({y}_{1}-{y}_{2})}^{2}}$}
\Note{Note that $(x_1 - x_2)^2 = (x_2 - x_1)^2$, but remember to be consistent or you will get the wrong answer.\\
$AB \neq \\ \sqrt{(x_1 - x_2)^2 + (y_2 - 2_1)^2}$.}

 

The following video provides a summary of the distance formula.
\setcounter{subfigure}{0}
\begin{figure}[H] % horizontal\label{m39107*uid99}
\textnormal{Khan academy video on distance formula}\vspace{.1in} \nopagebreak
\label{m39107*yt-media}\label{m39107*yt-video}
\raisebox{-5 pt}{ \includegraphics[width=0.5cm]{col11306.imgs/summary_www.png}} { (Video:  MG10108 )}
\vspace{2pt}
\vspace{.1in}
\end{figure}       

\begin{wex}{Distance formula I}{Find the distance between $S(-2;-5)$ and $Q(7;-2)$.}{
\westep{Draw a sketch}
% \begin{figure}
 \begin{center}

\scalebox{1} % Change this value to rescale the drawing.
{
\begin{pspicture}(0,-4.096719)(8.579062,4.1367188)
\rput(3.0,1.9032812){\psaxes[linewidth=0.028222222,arrowsize=0.05291667cm 2.0,arrowlength=1.4,arrowinset=0.4,labels=none,ticks=none,ticksize=0.10583333cm]{<->}(0,0)(-3,-6)(5,2)}
\psline[linewidth=0.028222222cm](1.8,-2.0367188)(6.3,0.9632813)
\psline[linewidth=0.028222222,linestyle=dashed,dash=0.17638889cm 0.10583334cm](1.8,-2.0367188)(6.3,-2.0367188)(6.3,0.9632813)
\psline[linewidth=0.028222222](6.06,-2.0367188)(6.06,-1.7832812)(6.3,-1.7832812)
\usefont{T1}{ptm}{m}{n}
\rput(2.7045312,1.6632812){$0$}
\usefont{T1}{ptm}{m}{n}
\rput(3.3545313,3.9332812){$y$}
\usefont{T1}{ptm}{m}{n}
\rput(8.204532,2.1432812){$x$}
\usefont{T1}{ptm}{m}{n}
\rput(5.744531,1.1932813){$T(7,-2)$}
\usefont{T1}{ptm}{m}{n}
\rput(2,-2.3367188){$S(-2,-5)$}
\psdots[dotsize=0.127](6.3,0.9632813)
\psdots[dotsize=0.127](1.8,-2.0367188)
\end{pspicture} 
}
 \end{center}
% \end{figure}

\westep{Assign values to $(x_1;y_1)$ and $(x_2;y_2)$}
Let the coordinates of $S$ be $(x_1;y_1)$ and the coordinates of $T$ be $(x_2;y_2)$.
\begin{equation*}
x_1 = -2 \hskip2em y_1 = -5 \hskip2em x_2 = -7 \hskip2em y_2 = -2
\end{equation*}
\westep{Write down the distance formula}
\begin{equation*}
d = \sqrt{(x_1 - x_2)^2 + (y_1 - y_2)^2}
\end{equation*}
\westep{Substitute values}
\begin{equation*}
\begin{array}{cl}
d_{\mbox{ST}} &= \sqrt{(2 - (-7))^2 + (-5- (-2))^2}\\
& = \sqrt{(-9)^2 + (-3)^2}\\
&= \sqrt{81 + 9}\\
&= \sqrt{90}\\
&= 9,5
\end{array}
\end{equation*}
\westep{Write the final answer}
The distance between $S$ and $T$ is $9,5$ units.
\vspace{2pt}
\vspace{.1in}
}
\end{wex}

\begin{wex}{Distance formula II}{Given $RS = 13$, $R(3;9)$ $S(8;y)$, find $y$.}{
\westep{Draw a sketch}
% \usepackage{pst-plot} % For axes
 \begin{center}
\scalebox{1} % Change this value to rescale the drawing.
{
\begin{pspicture}(0,-3.6167188)(7.8990626,3.6567187)
\rput(2.0,-1.6167188){\psaxes[linewidth=0.028222222,arrowsize=0.05291667cm 2.0,arrowlength=1.4,arrowinset=0.4,labels=none,ticks=none,ticksize=0.10583333cm]{<->}(0,0)(-2,-2)(5,5)}
\psline[linewidth=0.028222222,linestyle=dashed,dash=0.17638889cm 0.10583334cm](4.68,2.4567187)(3.1393335,-0.2630555)(4.64,-2.5032814)
\usefont{T1}{ptm}{m}{n}
\rput(1.7445313,-1.8767188){$0$}
\usefont{T1}{ptm}{m}{n}
\rput(2.3545313,3.4532812){$y$}
\usefont{T1}{ptm}{m}{n}
\rput(7.3,-1.4567188){$x$}
\usefont{T1}{ptm}{m}{n}
\rput(5.264531,2.7132812){$S_2(8;y_2)$}
\usefont{T1}{ptm}{m}{n}
\rput(5.0845313,-2.7567186){$S_1(8;y_1)$}
\psdots[dotsize=0.12](3.14,-0.22328125)
\psdots[dotsize=0.12](4.68,2.4367187)
\psdots[dotsize=0.12](4.66,-2.4632812)
\usefont{T1}{ptm}{m}{n}
\rput(2.6245313,-0.20671874){$R(3;9)$}
\psline[linewidth=0.04cm,linestyle=dotted,dotsep=0.16cm](4.68,3.4967186)(4.66,-3.5032814)
\end{pspicture} 
}
\end{center}
\westep{Assign values to $(x_1;y_1)$ and $(x_2;y_2)$:}
Let the coordinates of $R$ be $(x_1;y_1)$ and the coordinates of $S$ be $(x_2;y_2)$.
\begin{equation*}
x_1 = 3 \hskip2em y_1 = 9 \hskip2em x_2 = 8 \hskip2em y_2 = y
\end{equation*}
\westep{Write down the distance formula}
\begin{equation*}
d = \sqrt{(x_1 - x_2)^2 + (y_1 - y_2)^2}
\end{equation*}
\westep{Substitute values and solve for $y$}
\begin{equation*}
\begin{array}{cl}
13 &= \sqrt{(3 - 8)^2 + (9 - y)^2}\\
13^2 & = (-5)^2 + (81 - 18y + y^2)\\
0 &= y^2 - 18y - 63\\
&= (y+3) (y-21)\\
\therefore y &= -3 \mbox{ or } y = 21
\end{array}

\end{equation*}
\westep{Write the final answer}
$S$ is $(8;-3)$ or $(8;21)$
\vspace{2pt}
\vspace{.1in}
}
\end{wex}

\begin{exercises}{}{
\begin{enumerate}[label=\textbf{\arabic*}.]
\item Find the length of $AB$ if:
 \begin{enumerate}[noitemsep, label=\textbf{(\alph*)} ] 
\item $A(2;7)$ and $B(-3;5)$.
\item $A(-3;5)$ and $B(-9;1)$.
\item $A(x;y)$ and $B(x+4;y-1)$.
\end{enumerate}

\item The length of $CD=5$. Find the missing coordinate if:
 \begin{enumerate}[noitemsep, label=\textbf{(\alph*)} ] 
\item $C(6;-2)$ and $D(x;2)$.
\item $C(4;y)$ and $D(1;-1)$.
\end{enumerate}
\end{enumerate}

Find the answers with the shortcodes:}
\end{exercises}

%          \section{ Calculation of the gradient line}
%     \nopagebreak
%%%             \label{m39108} $ \hspace{-5pt}\begin{array}{cccccccccccc}   \includegraphics[width=0.75cm]{col11306.imgs/summary_video.png} &   \end{array} $ \hspace{2 pt}\raisebox{-5 pt}{} {(section shortcode: MG10109 )} \par 
%     
%     
%     

\section{Gradient between two points}
\Definition{Gradient}{The gradient between two points is determined by the ratio vertical change to horizontal change.}
Gradient $m$ describes the slope or steepness of the line joining two points. In the figure below, line $OQ$ is the least steep, and $OT$ is the steepest.

\setcounter{subfigure}{0}
\begin{figure}[H] % horizontal\label{m39107*id63458}
\begin{center}
\scalebox{.8}{
\begin{pspicture}(-5,-5)(5.5,5.5)
% \psaxes{<->}(0,0)(5,5)
% \psgrid[gridcolor=lightgray,linecolor=lightgray,subgriddiv=1,gridlabels=0.0cm](0,0)(-1,-1)(4,4)
\psaxes[linewidth=1pt,labels=none,ticks=none]{<->}(0,0)(-1,-1)(4,4)
\psline[linewidth=.05cm](0,0)(1,3.5)
\psline[linewidth=.05cm](0,0)(3,3.5)
\psline[linewidth=.05cm](0,0)(3,2)
\psline[linewidth=.05cm](0,0)(3,0.5)
\uput[ur](.9,3.5){\Large{$T$}}
% \uput[dl](-1,-2){\Large{$(x_{2},y_{2})$}}
\uput[ur](3,3.5){\Large{$S$}}
\uput[ur](3,2){\Large{$R$}}
\uput[ur](3,0.4){\Large{$Q$}}
\uput[dl](0,0){\Large{$O$}}
% \uput[ur](1.8,1){\Large{$(x_{1},y_{1})$}}
\uput[l](4.8,0){\Large{$x$}}
\uput[d](0,4.8){\Large{$y$}}
\end{pspicture}
}
\end{center}
\end{figure}        
To derive the formula for gradient, we consider and right-angled triangle foremd from $A~(x_1;y_1)$ and $B~(x_2;y_2)$ with hypotenuse $AB$, as shown in the diagram below. The gradient is determined by the ratio of the length of the vertical side of hte triangle to the length of the horizontal side of the triangle. The length of the vertical side of the triangle is the difference in $y$-values of points $A$ and $B$. The length of the horizontal side of the triangle is the difference in $x$-values of points $A$ and $B$. 
\setcounter{subfigure}{0}
\begin{figure}[H] % horizontal\label{m39107*id63458}
\begin{center}
\scalebox{.8}{
\begin{pspicture}(-5,-5)(5.5,5.5)
% \psaxes{<->}(0,0)(5,5)
% \psgrid[gridcolor=lightgray,linecolor=lightgray,subgriddiv=1,gridlabels=0.0cm](0,0)(-3,-3)(3,3)
\psaxes[linewidth=1pt,labels=none,ticks=none]{<->}(0,0)(-3,-3)(3,3)
\psline(-2,-2)(2,-2)(2,1)
\psline[linestyle=solid,linewidth=0.04cm](-2,-2)(2,1)
\uput[l](-2,-1.8){\Large{$A$}}
\uput[dl](-1,-2){\Large{$(x_{1},y_{1})$}}
\uput[dr](2,-2){\Large{$C$}}
\uput[u](2.3,0.5){\Large{$B$}}
\uput[ur](1.8,1){\Large{$(x_{2},y_{2})$}}
\uput[l](4,0){\Large{$x$}}
\uput[d](0,4){\Large{$y$}}
\uput[dl](0,0){\Large{$0$}}
\end{pspicture}
}
\end{center}
\end{figure}  
Therefore, gradient is determined using the following equation:

\Identity{Gradient}{Gradient ($m$) $= \dfrac{y_{2} - y_{1}}{x_{2} - x_{1}}$ or $= \dfrac{y_{1} - y_{2}}{x_{1} - x_{2}}$
 \\ [5pt] Remember to be consistent: $m \neq \dfrac{y_{1} - y_{2}}{x_{2} - x_{1}}$}
\Tip{Remember: we can also use $m$ for the gradient of a line}
\begin{wex}{Gradient between two points}{Find the gradient of the line between points E$(2;5)$ and F$(-3;9)$.}{
\westep{Draw a sketch}
\begin{center}
\scalebox{1} % Change this value to rescale the drawing.
{
\begin{pspicture}(0,-3.6267188)(6.6990623,3.6667187)
\rput(3.0,-1.6267188){\psaxes[linewidth=0.04,arrowsize=0.05291667cm 2.0,arrowlength=1.4,arrowinset=0.4,labels=none,ticks=none,ticksize=0.10583333cm]{<->}(0,0)(-3,-2)(3,5)}
\psdots[dotsize=0.12](4.82,0.69328123)
\psdots[dotsize=0.12](1.06,2.8732812)
\psline[linewidth=0.04cm](4.84,0.67328125)(1.08,2.8532813)
\usefont{T1}{ptm}{m}{n}
\rput(6.324531,-1.3767188){$x$}
\usefont{T1}{ptm}{m}{n}
\rput(3.3045313,3.4632812){$y$}
\usefont{T1}{ptm}{m}{n}
\rput(5.094531,0.5232813){$E~(2;5)$}
\usefont{T1}{ptm}{m}{n}
\rput(1.2945312,3.0832813){$F~(-3;9)$}
\usefont{T1}{ptm}{m}{n}
\rput(2.7445312,-1.8367188){$0$}
\end{pspicture} 
}
\end{center}
\westep{Assign values to $(x_1;y_1)$ and $(x_2;y_2)$}
Let the coordinates of $E$ be $(x_1;y_1)$ and the coordinates of $F$ be $(x_2;y_2)$.
\begin{equation*}
x_1 = 2 \hskip2em y_1 = 5 \hskip2em x_2 = -3 \hskip2em y_2 = 9
\end{equation*}
\westep{Write down the formula for gradient}
\begin{equation*}
m = \frac{y_2 - y_1}{x_2 - x_1}
\end{equation*}
\westep{Substitute known values}
\begin{equation*}
\begin{array}{cl}
m_{EF} &= \frac{9 - 5}{-3 - 2}\\
&= \frac{4}{-5}
\end{array}
\end{equation*}
\westep{Write the final answer}
The gradient of $EF = \dfrac{4}{-5}$

}
\end{wex}


\begin{wex}{Gradient between two points}{Given $G(7;-9)$ and $H(x;0)$, with $m_{GH}= 3$, Find $x$.}{
\westep{Draw a sketch}
\begin{center}
\scalebox{1} % Change this value to rescale the drawing.
{
\begin{pspicture}(0,-3.6867187)(7.1790624,3.7267187)
\rput(3.0,1.3132813){\psaxes[linewidth=0.04,arrowsize=0.05291667cm 2.0,arrowlength=1.4,arrowinset=0.4,labels=none,ticks=none,ticksize=0.10583333cm]{<->}(0,0)(-3,-3)(3,2)}
\psdots[dotsize=0.12](4.9,1.2932812)
\psdots[dotsize=0.12](3.9,-0.74671876)
\psline[linewidth=0.04cm](3.88,-0.76671875)(4.9,1.2732812)
\usefont{T1}{ptm}{m}{n}
\rput(6.324531,1.5632813){$x$}
\usefont{T1}{ptm}{m}{n}
\rput(3.1445312,3.5232813){$y$}
\usefont{T1}{ptm}{m}{n}
\rput(5.324531,1.6032813){$H~(x;0)$}
\usefont{T1}{ptm}{m}{n}
\rput(4.134531,-0.99671876){$G~(7;-9)$}
\usefont{T1}{ptm}{m}{n}
\rput(2.7445312,1.0432812){$0$}
\usefont{T1}{ptm}{m}{n}
\rput(5.7745314,0.28984374){$m_{GH} = 3$}
\end{pspicture} 
}
\end{center}
\westep{Assign values to $(x_1;y_1)$ and $(x_2;y_2)$}
Let the coordinates of $G$ be $(x_1;y_1)$ and the coordinates of $H$ be $(x_2;y_2)$.
\begin{equation*}
x_1 = 7 \hskip2em y_1 = -9 \hskip2em x_2 = x \hskip2em y_2 = 0
\end{equation*}
\westep{Write down the formula for gradient}
\begin{equation*}
m = \frac{y_2 - y_1}{x_2 - x_1}
\end{equation*}
\westep{Substitute values and solve for $x$}
\begin{equation*}
\begin{array}{cl}
3 &= \frac{0 - (-9)}{x - 7}\\
3(x-7)&= 9\\
x-7 &= \frac{9}{3}\\
x-7 &= 3\\
x &= 3 + 7\\
&= 10 \\
\end{array}
\end{equation*}
\westep{Write the final answer}
The coordinates of $H$ are $(10;0)$.
\vspace{2pt}
\vspace{.1in}
}
\end{wex}


\begin{exercises}{}
 \begin{enumerate}[noitemsep, label=\textbf{(\alph*)} ]
\item Find the gradient of $AB$ if:
 \begin{enumerate}[noitemsep, label=\textbf{(\alph*)} ] 
\item $A(7;10)$ and $B(-4;1)$.
\item $A(-5;-9)$ and $B(3;2)$.
\item $A(x-3;y)$ and $B(x;y+4)$.
\end{enumerate}

\item If the gradient of $CD=\frac{2}{3}$, find $P$:
\begin{enumerate} \begin{enumerate}[noitemsep, label=\textbf{(\alph*)} ] 
\item $C(16;2)$ and $D(8;P)$.
\item $C(3;2P)$ and $D(9;14)$.
\end{enumerate}
\end{enumerate}

Find the answers with the shortcodes:
}
\end{exercises}
\subsection*{Straight line}
\Definition{Straight line}{A straight line is a set of points with a constant gradient between any two of the points.}
Consider the diagram below with points $A(x;y)$, $B(x_2;y_2)$ and $C(x_1;y_1)$. 
\begin{center}
 \scalebox{1} % Change this value to rescale the drawing.
{
\begin{pspicture}(0,-3.0667188)(7.5290623,3.1067188)
\rput(3.97,-0.06671875){\psaxes[linewidth=0.04,arrowsize=0.05291667cm 2.0,arrowlength=1.4,arrowinset=0.4,labels=none,ticks=none,ticksize=0.10583333cm]{<->}(0,0)(-3,-3)(3,3)}
\psdots[dotsize=0.12](3.39,0.47328126)
\psdots[dotsize=0.12](4.57,2.2132812)
\psline[linewidth=0.04cm](2.17,-1.3467188)(4.69,2.3932812)
\usefont{T1}{ptm}{m}{n}
\rput(7.1545315,0.14328125){$x$}
\usefont{T1}{ptm}{m}{n}
\rput(4.2545314,2.9032812){$y$}
\usefont{T1}{ptm}{m}{n}
\rput(5.514531,2.1832812){$A(x;y)$}
\usefont{T1}{ptm}{m}{n}
\rput(3.6745312,-0.33671874){$0$}
\psdots[dotsize=0.12](2.51,-0.86671877)
\usefont{T1}{ptm}{m}{n}
\rput(2.5845313,0.5832813){$B(x_2;y_2)$}
\usefont{T1}{ptm}{m}{n}
\rput(1.7,-0.7967188){$C(x_1;y_1)$}
\end{pspicture} 
}
\end{center}
We have $m_{AB} = m_{BC}=m_{AC}$ \par
and $m = \dfrac{y_2-y_1}{x_2-x_1} = \dfrac{y_1-y_2}{x_1-x_2}$\par

The general formula for a straight line is $\dfrac{y-y_1}{x-x_1} = \dfrac{y_2-y_1}{x_2-x_1}$ (where $(x;y)$ is any point on the line)\par

The standard form of the straight line equation is $y=mx+c$ where $m$ is the gradient and $c$ is the $y$-intercept.

\begin{wex}{Finding the equation of a straight line}
 {Find the equation of the straight line through $P(-1;-5)$ and $Q(5;4)$.}
{
\westep{Draw a sketch}
\begin{center}
\scalebox{1} % Change this value to rescale the drawing.
{
\begin{pspicture}(0,-3.0667188)(6.5590625,3.1067188)
\rput(3.0,-0.06671875){\psaxes[linewidth=0.04,arrowsize=0.05291667cm 2.0,arrowlength=1.4,arrowinset=0.4,labels=none,ticks=none,ticksize=0.10583333cm]{<->}(0,0)(-3,-3)(3,3)}
\psdots[dotsize=0.12,dotangle=-5.9493704](4.5090275,2.306018)
\psline[linewidth=0.04cm](2.14,-0.52671874)(4.48,2.3132813)
\usefont{T1}{ptm}{m}{n}
\rput(6.184531,0.14328125){$x$}
\usefont{T1}{ptm}{m}{n}
\rput(3.2845314,2.9032812){$y$}
\usefont{T1}{ptm}{m}{n}
\rput(4.8545313,2.6232812){$Q(5;4)$}
\usefont{T1}{ptm}{m}{n}
\rput(2.7045312,-0.33671874){$0$}
\psdots[dotsize=0.12,dotangle=-5.9493704](2.1408823,-0.5438744)
\usefont{T1}{ptm}{m}{n}
\rput(1.7045312,-0.81671876){$P(-1;-5)$}
\end{pspicture} 
}
\end{center}
\westep{Assign values}
Let $(x;y)$ be any point on the line. \\
$x_1 = 5 \hskip2em y_1 = 4 \hskip2em x_1 = x \hskip2em y_2 = -5$


\westep{Write down the general formula}
\begin{align*}
\dfrac{y-y_1}{x-x_1} &= \dfrac{y_2-y_1}{x_2-x_1}
\end{align*}
\westep{Substitute values and make $y$ the subject of the equation}
\begin{align*}
 \dfrac{y-(-5)}{x-(-1)} &= \dfrac{4-(-5)}{5-(-1)} \\

 \dfrac{y+5}{x+1)} &= \dfrac{9}{6} = \dfrac{3}{2}\\
2(y+5) &=3(x+1)\\
2y +10&=3x+3\\
2y&=3x-7\\
y&=\frac{3}{2}x - \frac{7}{2}
\end{align*}
\westep{Write the final answer}
The equation of the straight line is $y&=\frac{3}{2}x - \frac{7}{2}$.
}


\end{wex}

\subsection*{Parallel and perpendicular lines}    
%         \label{m39108*eip-332}We can use the gradient of a line to determine if two lines are parallel or perpendicular. If the lines are parallel (Figure~\ref{fig:parallelperpendicular}a) then they will have the same gradient, i.e. ${m}_{\mbox{AB}}={m}_{\mbox{CD}}$. If the lines are perpendicular (Figure~\ref{fig:parallelperpendicular}b) than we have: 
%     \setcounter{subfigure}{0}
%  	\begin{figure}[H] % horizontal\label{m39107*id63458}
%     \begin{center}
% \scalebox{.8}{
% \begin{pspicture}(-5,-5)(5.5,5.5)
% % \psaxes{<->}(0,0)(5,5)
% \rput(-2,-2){
% \psline[linewidth=.05cm](0,0)(0,3)
% \psline[linewidth=.05cm](1,0)(1,3)
% \uput[ur](-1,2.8){\Large{$a)$}}
% \uput[d](0,0){\Large{$A$}}
% \uput[u](0,3){\Large{$B$}}
% \uput[d](1,0){\Large{$C$}}
% \uput[u](1,3){\Large{$D$}}}
% 
% \rput(3.4,2){
% \psline[linewidth=.05cm](1,-4)(-2,-1)
% \psline[linewidth=.05cm](-2,-4)(1,-1)
% \uput[ur](-3.2,-1.2){\Large{$b)$}}
% \uput[dr](1,-4){\Large{$A$}}
% \uput[ul](-2,-1){\Large{$B$}}
% \uput[dl](-2,-4){\Large{$C$}}
% \uput[ur](1,-1){\Large{$D$}}}
% \end{pspicture}
% }
%     \end{center}
% \caption{a) Parallel and b) perpendicular lines}
% \label{fig:parallelperpendicular}
%  \end{figure}       
Two lines have that run parallel to each other have equal gradients. \\
In other words: $\mbox{gradient}_{AB}=\mbox{gradient}_{CD}$. \par
If two lines intersect perpendicularly then the product of their gradients is equal to $-1$. \\

If line $WX \perp $ line $ YZ$ then $m_{WX} \times m_{YZ} = -1$.

\begin{wex}{Parallel lines}{Prove that the line $AB$ with points $A(2;4)$ and $B(8,16)$ is parallel to the line $y = 2x-2$.}{
\westep{Draw a sketch}

\begin{center}
 \scalebox{1} % Change this value to rescale the drawing.
{
\begin{pspicture}(0,-3.0667188)(6.5590625,3.1067188)
\rput(3.0,-0.06671875){\psaxes[linewidth=0.04,arrowsize=0.05291667cm 2.0,arrowlength=1.4,arrowinset=0.4,labels=none,ticks=none,ticksize=0.10583333cm]{<->}(0,0)(-3,-3)(3,3)}
\psdots[dotsize=0.12,dotangle=-5.9493704](4.5090275,2.306018)
\psline[linewidth=0.04cm](1.56,-1.7267188)(4.52,2.3132813)
\usefont{T1}{ptm}{m}{n}
\rput(6.184531,0.14328125){$x$}
\usefont{T1}{ptm}{m}{n}
\rput(3.2845314,2.9032812){$y$}
\usefont{T1}{ptm}{m}{n}
\rput(4.9445314,2.6232812){$B(8;16)$}
\usefont{T1}{ptm}{m}{n}
\rput(2.7845314,-0.31671876){$0$}
\usefont{T1}{ptm}{m}{n}
\rput(2.8345313,0.94328123){$A(2;4)$}
\psline[linewidth=0.04cm](1.9,-2.0667188)(4.86,1.9732813)
\psdots[dotsize=0.12](3.4,0.79328126)
\usefont{T1}{ptm}{m}{n}
\rput(3.1945312,-0.69671875){$-2$}
\rput(5.2145314,1.2232813){$y=2x-2$}
\end{pspicture} 
}

\end{center}
Be careful - some lines may look parallel but are not!

\westep{Write down the formula for gradient}
\begin{equation*}
m = \dfrac{y_2-y_1}{x_2-x_1}
\end{equation*}
\westep{Substitute values to find the gradient for line $AB$}
\begin{equation*}
\begin{array}{cl}
m_{AB} &= \dfrac{16 - 4}{8 - 2}\\[5pt]
&= \dfrac{12}{6}\\
&= 2
\end{array}
\end{equation*}
\westep{Check that $y$ is in the standard form $y=mx+c$}
\begin{equation*}
\begin{array}{cl}
y&=2x-2\\
\therefore m_{y}&= 2
\end{array}
\end{equation*}
\westep{Write the final answer}
\begin{equation*}
\begin{array}{cl}
m_{AB} &= m_{y}\\

\end{array}
\end{equation*}
therefore line $AB$ is parallel to $y=2x-2$.
}
\end{wex}

Perpendicular lines have gradients that are the negative inverse of each other. This means that when you multiply the gradients together, the answer will be -1.: $-\frac{1}{m_{AB}}=m_{CD}$

\begin{wex}{Perpendicular lines}{ Line $AB$ is perpendicular to line $CD$. Find $y$ given $A(2;-3)$, $B(-2;6)$, $C(4;3)$ and $D(7;y)$.}{
\westep{Draw a sketch}
\begin{center}
 \scalebox{1} % Change this value to rescale the drawing.
{
\begin{pspicture}(0,-3.0667188)(6.7390623,3.1067188)
\rput(3.0,-0.06671875){\psaxes[linewidth=0.04,arrowsize=0.05291667cm 2.0,arrowlength=1.4,arrowinset=0.4,labels=none,ticks=none,ticksize=0.10583333cm]{<->}(0,0)(-3,-3)(3,3)}
\psdots[dotsize=0.12,dotangle=-5.9493704](4.089028,0.96601814)
\psline[linewidth=0.04cm](1.62,-1.5267187)(5.68,2.5732813)
\usefont{T1}{ptm}{m}{n}
\rput(6.184531,0.14328125){$x$}
\usefont{T1}{ptm}{m}{n}
\rput(3.2845314,2.9032812){$y$}
\usefont{T1}{ptm}{m}{n}
\rput(1.3445313,2.4432812){$B(8;16)$}
\usefont{T1}{ptm}{m}{n}
\rput(2.6845312,-0.25671875){$0$}
\usefont{T1}{ptm}{m}{n}
\rput(4.264531,-0.7967188){$A(2;-3)$}
\psline[linewidth=0.04cm](1.52,2.1732812)(4.2,-0.56671876)
\psdots[dotsize=0.12](1.56,2.1732812)
\psdots[dotsize=0.12](4.18,-0.52671874)
\usefont{T1}{ptm}{m}{n}
\rput(3.9845312,1.2232813){$C(4;3)$}
\psline[linewidth=0.04cm,linestyle=dashed,dash=0.16cm 0.16cm](5.1,2.8132813)(5.12,-0.44671875)
\usefont{T1}{ptm}{m}{n}
\rput(5.914531,1.9032812){$D(7;y)$}
\psdots[dotsize=0.12](5.12,2.0332813)
\end{pspicture} 
}
\end{center}


\westep{Write the formulaa for the product of perpendicular gradients}
\begin{equation*}
m_{AB} \times m_{CD} = -1
\end{equation*}
\westep{Substitute values and solve for $y$}
\begin{equation*}
\begin{array}{rl}
\frac{6 - (-3)}{-2 -2} \times \frac{y - 3}{7 - 4} &= -1\\
\frac{9}{-4} \times \frac{y-3}{3} &= -1\\
\frac{y-3}{3} &= -1 \times \frac{-4}{9}\\
\frac{y-3}{3} &= \frac{4}{9}\\
y-3 &= \frac{4}{9} \times 3\\
y-3 &= \frac{4}{3}\\
y &= \frac{4}{3} + 3\\
&= \frac{4 + 9}{3}\\
&= \frac{13}{3}\\
&= 4 \frac{1}{3}
\end{array}
\end{equation*}
\westep{Write the final answer}
Therefore the coordinates of $D$ are $(7; 4\frac{1}{3})$.
}
\end{wex}

\subsection*{Horizontal and vertical lines}

A line that runs parallel to the $x$-axis is called a horizontal line and has a gradient of zero. This is
because there is no vertical change in the gradient of the line:\par
$m = \dfrac{\mbox{change in }y}{\mbox{change in }x} = \dfrac{0}{\mbox{change in }x} =0$\par

A line runs parallel to the $y$-axis,is called a vertical line and its gradient is undefined. This is because there is no horizontal change in the gradient of the line:\par
$m = \dfrac{\mbox{change in }y}{\mbox{change in }x} = \dfrac{\mbox{change in }y}{0}=$ undefined\par

\subsection*{Collinear points}

Points are collinear if they lie on the same line. There are two methods to prove that points are collinear; the gradient method, and a
longer method using the distance formula.

If three points are collinear then the line joining them will have the same gradient at any point along the line.
Therefore to prove collinearity, we must prove that two of the gradients between any of the three points are
equal.

\begin{wex}{Collinear points}{Prove that $A(-3;3)$, $B(0;5)$ and $C(3;7)$ are collinear.}{
\westep{Find any two gradients:}
\begin{equation*}
m_{AB} = \frac{5-3}{0-(-3)} = \frac{2}{3}
\end{equation*}
\begin{equation*}
m_{BC} = \frac{7-5}{3-0} = \frac{2}{3}
\end{equation*}
\westep{Establish if the points are collinear and explain your answer:}
\begin{equation*}
m_{AB} = m_{BC}
\end{equation*}
$\therefore$ Points $A$, $B$ and $C$ are collinear.
}
\end{wex}

To prove that three points are collinear using the distance formula requires finding the
distances between each pair points and then proving that the sum of the two smaller distances
equals the longest distance.

To show this, let's do the same worked example.

\begin{wex}{Collinear points}{Prove that $A(-3;3)$, $B(0;5)$ and $C(3;7)$ are collinear.}{
\westep{Find all three distances:}
\begin{equation*}
d_{AB} = \sqrt{(-3 - 0)^2 + (3 - 5)^2} = \sqrt{(-3)^2 + (-2)^2} = \sqrt{9 + 4} = \sqrt{13}
\end{equation*}
\begin{equation*}
d_{BC} = \sqrt{(0 - 3)^2 + (5 - 7)^2} = \sqrt{(-3)^2 + (-2)^2} = \sqrt{9 + 4} = \sqrt{13}
\end{equation*}
\begin{equation*}
d_{AC} = \sqrt{((-3) - 3)^2 + (3 - 7)^2} = \sqrt{(-6)^2 + (-4)^2} = \sqrt{36 + 16} = \sqrt{52}
\end{equation*}
\westep{Find the sum of the two shorter distances:}
\begin{equation*}
d_{AB} + d_{BC} = \sqrt{13} + \sqrt{13} = 2\sqrt{13} = \sqrt{4 \times 13} = \sqrt{52}
\end{equation*}
\westep{Prove that the points are collinear:}
\begin{equation*}
d_{AB} + d_{BC} = d_{AC}
\end{equation*}
$\therefore$ Points $A$, $B$ and $C$ are collinear.
}
\end{wex}

\begin{exercises}{}
\begin{enumerate}
\item Determine whether $AB$ is parallel, perpendicular or neither to $CD$ if:
\begin{enumerate}
\item $A(3;-4)$, $B(5;2)$, $C(-1;-1)$, $D(7;23)$.
\item $A(3;-4)$, $B(5;2)$, $C(-1;-1)$, $D(0;-4)$.
\item $A(3;-4)$, $B(5;2)$, $C(-1;-1)$, $D(1;2)$.
\end{enumerate}

\item Determine whether the following points are collinear:
\begin{enumerate}
\item $E(0;3)$, $F(-2;5)$, $G(2;1)$.
\item $H(-3;-5)$, $I(-0;0)$, $J(6;10)$.
\item $K(-6;2)$, $L(-3;1)$, $M(1;-1)$.
\end{enumerate}
\item Points $P(-6;2)$, $Q(2;-2)$ and $R(-3;r)$ lie on a straight line. Find the value of $r$.
\end{enumerate}
\end{exercises}

The following video provides a summary of the gradient of a line.
\setcounter{subfigure}{0}
\begin{figure}[H] % horizontal\label{m39108*uid993}
\textnormal{Gradient of a line}\vspace{.1in} \nopagebreak
\label{m39108*yt-media1}\label{m39108*yt-video1}
\raisebox{-5 pt}{ \includegraphics[width=0.5cm]{col11306.imgs/summary_www.png}} { (Video:  MG10110 )}
\vspace{2pt}
\vspace{.1in}
\end{figure}      
%          \section{ Midpoint of a line}
%     \nopagebreak
%%%             \label{m39119} $ \hspace{-5pt}\begin{array}{cccccccccccc}   \includegraphics[width=0.75cm]{col11306.imgs/summary_video.png} &   \end{array} $ \hspace{2 pt}\raisebox{-5 pt}{} {(section shortcode: MG10111 )} \par 
%     
%     
%     
\section{Straight line}

\section{Mid-point of a line}
\nopagebreak
Sometimes, knowing the coordinates of the middle point or \textsl{mid-point} of a line is useful. For example, what is the mid-point of the line between point $P$ with coordinates $(2;1)$ and point $Q$ with coordinates $(-2;-2)$.\par 
The coordinates of the mid-point of any line between any two points $A$ and $B$ with coordinates $({x}_{1};{y}_{1})$ and $({x}_{2};{y}_{2})$, is generally calculated as follows. Let the mid-point of $AB$ be at point $S$ with coordinates $(X;Y)$. The aim is to calculate $X$ and $Y$ in terms of $({x}_{1};{y}_{1})$ and $({x}_{2};{y}_{2})$.\par 
\setcounter{subfigure}{0}
\begin{figure}[H] % horizontal\label{m39107*id63458}
\begin{center}
\scalebox{.8}{
\begin{pspicture}(-5,-5)(5.5,5.5)
\psgrid[gridcolor=lightgray,linecolor=lightgray,subgriddiv=1,gridlabels=0.0cm](0,0)(0,0)(5,5)
\psaxes[linewidth=1pt,labels=none,ticks=none]{->}(0,0)(0,0)(5,5)
\psline[linewidth=.05cm](0,0)(5,5)
\uput[dl](0,0){\Large{$A (x_{1},y_{1})$}}
\uput[u](0,-.2){\qdisk(0,0){3pt}}
\uput[r](2.6,2.4){\Large{$S (X,Y)$}}
\uput[u](2.5,2.3){\qdisk(0,0){3pt}}
\uput[ur](5,5){\Large{$B (x_{2},y_{2})$}}
\uput[u](5,4.8){\qdisk(0,0){3pt}}
\uput[l](5.8,0){\Large{$x$}}
\uput[d](0,5.8){\Large{$y$}}
\end{pspicture}
}
\end{center}
\end{figure}      

\begin{eqnarray*}
X & = & \frac{{x}_{1} + {x}_{2}}{2} \\ 
Y & = & \frac{{y}_{1} + {y}_{2}}{2} \\  
\end{eqnarray*}
From this we obtain the mid-point of a line:
\Identity{Midpoint}{$S(\frac{{x}_{1} + {x}_{2}}{2};\frac{{y}_{1}+{y}_{2}}{2}) $}
\label{m39119*id68788}Then the coordinates of the mid-point ($S$) of the line between point $P$ with coordinates $(2;1)$ and point $Q$ with coordinates $(-2;-2)$ is:


\begin{eqnarray*}
X & = & \frac{{x}_{1} + {x}_{2}}{2} \\ 
& = & \frac{-2 + 2}{2} \\ 
& = & 0 \\ 
Y & = & \frac{{y}_{1} + {y}_{2}}{2} \\ 
& = & \frac{-2 + 1}{2} \\ 
& = & -\frac{1}{2} 
\end{eqnarray*}
The mid-point is at $S(0;-\frac{1}{2})$. 
It can be confirmed that the distance from each end point to the mid-point is equal. The coordinate of the mid-point $S$ is $(0;-0,5)$.\par 

\begin{eqnarray*}
PS & = & \sqrt{{({x}_{1} - {x}_{2})}^{2} + {({y}_{1} - {y}_{2})}^{2}} \\ 
& = & \sqrt{{(0 - 2)}^{2} + {(-0,5 - 1)}^{2}} \\ 
& = & \sqrt{{(-2)}^{2} + {(-1,5)}^{2}} \\ 
& = & \sqrt{4 + 2,25} \\ 
& = & \sqrt{6,25}
\end{eqnarray*}
and
\begin{eqnarray*}
QS & = & \sqrt{{({x}_{1} - {x}_{2})}^{2} + {({y}_{1} - {y}_{2})}^{2}} \\ 
& = & \sqrt{{(0 - (-2))}^{2} + {(-0,5 - (-2))}^{2}} \\ 
& = & \sqrt{{(0 + 2)}^{2}{+(-0,5 + 2)}^{2}} \\ 
& = & \sqrt{{(2)}^{2}{+(-1,5)}^{2}} \\ 
& = & \sqrt{4 + 2,25} \\ 
& = & \sqrt{6,25}
\end{eqnarray*}
\label{m39119*id69633}It can be seen that $PS=QS$ as expected.\par 
\setcounter{subfigure}{0}
\begin{figure}[H] % horizontal\label{m39107*id63458}
\begin{center}
\scalebox{.8}{
\begin{pspicture}(-5,-5)(5.5,5.5)
% \psaxes{<->}(0,0)(5,5)
\psgrid[gridcolor=lightgray,linecolor=lightgray,subgriddiv=1,gridlabels=0.0cm](0,0)(-3,-3)(3,3)
\psaxes[linewidth=1pt,labels=none,ticks=none]{->}(0,0)(-3,-3)(3,3)
\psline[linewidth=.05cm](-2,-2)(2,1)
% \psline[linewidth=.05cm](0,0)(3,3.5)
% \psline[linewidth=.05cm](0,0)(3,2)
% \psline[linewidth=.05cm](0,0)(3,0.5)
% \uput[ur](.9,3.5){\Large{$T$}}
\uput[dl](-1,-2){\Large{$Q (-2,-2)$}}
\uput[u](-2,-2.2){\qdisk(0,0){3pt}}
\uput[r](0,-.7){\Large{$S$}}
\uput[r](0,-1.2){\Large{mid-point}}
\uput[u](0,-.7){\qdisk(0,0){3pt}}
\uput[ur](2,1){\Large{$P (2,1)$}}
\uput[u](2,.8){\qdisk(0,0){3pt}}
% \uput[l](4.8,0){\Large{$x$}}
% \uput[d](0,4.8){\Large{$y$}}
\end{pspicture}
}
\end{center}
\end{figure} 
The following video provides a summary of the mid-point of a line.
\setcounter{subfigure}{0}
\begin{figure}[H] % horizontal\label{m39119*uid666}
\textnormal{Khan academy video on mid-point of a line}\vspace{.1in} \nopagebreak
\label{m39119*yt-media2}\label{m39119*yt-video2}
\raisebox{-5 pt}{ \includegraphics[width=0.5cm]{col11306.imgs/summary_www.png}} { (Video:  MG10112 )}
\vspace{2pt}
\vspace{.1in}
\end{figure}      
%          \section{ Summary \& Exercises}
%     \nopagebreak
%%%             \label{m39167} $ \hspace{-5pt}\begin{array}{cccccccccccc}   \end{array} $ \hspace{2 pt}\raisebox{-5 pt}{\includegraphics[width=0.5cm]{col11306.imgs/summary_www.png}} {(section shortcode: MG10113 )} \par 
%     
%     
% 
\begin{wex}{Midpoint formula I}{Find the mid-point of line $AB$ if $A$ is at $(6;2)$ and $B$ is at $(-5;-1)$.}{
\westep{Assign values to $(x_1;y_1)$ and $(x_2;y_2)$:}
\begin{equation*}
x_1= 6 \hskip2em y_1=2 \hskip2em x_2=-5 \hskip y_2=-1
\end{equation*}
\westep{Write down the formula you are using:}
\begin{equation*}
(\frac{x_1+x_2}{2};\frac{y_1+y_2}{2})
\end{equation*}
\westep{Substitute known variables and solve:}
\begin{equation*}
(\frac{6-5}{2};\frac{2-1}{2}) = (\frac{1}{2};\frac{1}{2})
\end{equation*}
\westep{Explain your answer:}
$(\frac{1}{2};\frac{1}{2})$ is the mid-point of $AB$.
}
\end{wex}

\begin{wex}{Midpoint formula II}{The line joining $C(-2;4)$ and $D(x;y)$ has the mid-point $M(1;-3)$. Find point $D$.}{
\westep{Assign values to $(x_1;y_1)$ and $(x_2;y_2)$:}
Let the coordinates of $C$ be $(x_1;y_1)$ and the co-orindates of $D$ be $(x_2;y_2)$.
\begin{equation*}
x_1=-2 \hskip2em y_1=4 \hskip2em x_2=x \hskip2em y_2=y
\end{equation*}
\westep{Write down the formula you are using:}
\begin{equation*}
(\frac{x_1+x_2}{2}; \frac{y_1+y_2}{2})
\end{equation*}
\westep{Substitue known variables and solve:}
\begin{equation*}
\begin{array}{rllrl}
1&=\frac{-2+x}{2}&\hskip10em&-3&=\frac{4+y}{2}\\
1\times 2&=-2+x&&-3\times 2&=4+y\\
2&=-2+x&&-6&=4+y\\
x&=2+2&&y&=-6-4\\
x&=4&&y&=-10\\
\end{array}
\end{equation*}
\westep{Explain your answer:}
Point $D$ is at $(4;-10)$.
}
\end{wex}

\begin{wex}{Midpoint formula III}{Points $E(-1;0)$ , $F(0;3)$ , $G(8;11)$ and $H(x;y)$ are points on the Cartesian plane. Find $H(x;y)$ if $EFGH$ is a parallelogram.}{
Remember, the diagonals of a parallelogram bisect each other. This mean that the mid-point of $EG$
will be equal to the mid-point of $FH$. If we find the first mid-point, we can use it to find point $H$.
\westep{For line $EG$, assign values to $(x_1;y_1)$ and $(x_2;y_2)$:}
Let the coordinates of $E$ be $(x_1;y_1)$ and the co-orindates of $G$ be $(x_2;y_2)$.
\begin{equation*}
x_1=-1 \hskip2em y_1=0 \hskip2em x_2=8 \hskip2em y_2=11
\end{equation*}
\westep{Write down the formula you are using:}
\begin{equation*}
(\frac{x_1+x_2}{2}; \frac{y_1+y_2}{2})
\end{equation*}
\westep{Substitue known variables and solve:}
\begin{equation*}
(\frac{-1+8}{2}; \frac{0+11}{2}) = (\frac{7}{2};\frac{11}{2})
\end{equation*}
\westep{For line $FH$, assign new values to $(x_1;y_1)$ and $(x_2;y_2)$:}
Let the coordinates of $F$ be $(x_1;y_1)$ and the co-orindates of $H$ be $(x_2;y_2)$.
\begin{equation*}
x_1=0 \hskip2em y_1=3 \hskip2em x_2=x \hskip2em y_2=y
\end{equation*}
\westep{Use the mid-point to find H:}
\begin{equation*}
\begin{array}{rllrl}
\frac{7}{2}&=\frac{0+x}{2}&\hskip10em&\frac{11}{2}&=\frac{3+y}{2}\\
7&=x+0&&11&=3+y\\
x&=7&&y&=8\\
\end{array}
\end{equation*}
\westep{Explain your answer:}
Point $H$ is at $(7;8)$.
}
\end{wex}

\begin{exercises}{Midpoint Formula}
\begin{enumerate}
\item Find the mid-points of the following lines:
\begin{enumerate}
\item $A(2;5)$, $B(-4;7)$.
\item $C(5;9)$, $D(23;55)$.
\item $E(x+2;y-1)$, $F(x-5;y-4)$.
\end{enumerate}

\item The mid-point $M$ of $PQ$ is (3;9). Find $P$ is $Q$ is $(-2;5)$:

\item $MNOP$ is a parallelogram with the points $M(5;3)$, $N(2;1)$ and $O(7;-3)$. Find point $P$.
\end{enumerate}
\end{exercises}    

\summary
\begin{itemize}[noitemsep]
\item Figures can be represented on the Cartesian plane\item The formula for finding the distance between two points is: 
\begin{equation}
\mathsf{Distance}=\sqrt{{({x}_{1}-{x}_{2})}^{2}+{({y}_{1}-{y}_{2})}^{2}}
\end{equation}
\item The formula for finding the gradient of a line is: 
\begin{equation}
\mathsf{Gradient}=\frac{{y}_{2}-{y}_{1}}{{x}_{2}-{x}_{1}}
\end{equation}
\item The formula for finding the mid-point between two points is: 
\begin{equation}
S(\frac{{x}_{1}+{x}_{2}}{2};\frac{{y}_{1}+{y}_{2}}{2})
\end{equation}
\item  If two lines are parallel then they will have the same gradient, i.e. ${m}_{\mathsf{AB}}={m}_{\mathsf{CD}}$. If two lines are perpendicular than we have: $-\frac{1}{{m}_{\mathsf{AB}}}={m}_{\mathsf{CD}}$\end{itemize}
\par 

\begin{eocexercises}{}
\begin{enumerate}[noitemsep, label=\textbf{\arabic*}. ] 
\item 
Represent the following figures on the Cartesian plane: 
\begin{enumerate}[noitemsep, label=\textbf{\alph*}. ] 
\item Triangle $DEF$ with $D(1;2)$, $E(3;2)$ and $F(2;4)$ 
\item Quadrilateral $GHIJ$ with $G(2;-1)$, $H(0;2)$, $I(-2;-2)$ and $J(1;-3)$
\item Quadrilateral $MNOP$ with $M(1;1)$, $N(-1;3)$, $O(-2;3)$ and $P(-4;1)$ 
\item Quadrilateral $WXYZ$ with $W(1;-2)$, $X(-1;-3)$, $Y(2;-4)$ and $Z(3;-2)$
\end{enumerate}
\item 
In the diagram given the vertices of a quadrilateral are $F(2;0)$, $G(1;5)$, $H(3;7)$ and $I(7;2)$.
\setcounter{subfigure}{0}
\begin{figure}[H] % horizontal\label{m39107*id63458}
\begin{center}
\scalebox{.8}{
\begin{pspicture}(-5,-5)(5.5,5.5)
% \psaxes{<->}(0,0)(5,5)
\psgrid[gridcolor=lightgray,linecolor=lightgray,subgriddiv=1](0,0)(-1,-1)(7,7)
\psaxes[linewidth=1pt,labels=none,ticks=none]{->}(0,0)(-1,-1)(7,7)
\pspolygon[linewidth=.05cm,fillstyle=solid,fillcolor=lightgray](2,0)(1,5)(3,7)(7,2)(2,0)
% \psline[linewidth=.05cm](0,0)(3,3.5)
% \psline[linewidth=.05cm](0,0)(3,2)
% \psline[linewidth=.05cm](0,0)(3,0.5)
% \uput[ur](.9,3.5){\Large{$T$}}
\uput[d](2,-.2){\normalsize{$F (2,0)$}}
\uput[ul](1.2,5){\normalsize{$G (1,5)$}}
\uput[u](3,7){\normalsize{$H (3,7)$}}
\uput[r](7,2){\normalsize{$I (7,2)$}}
\uput[l](8,0){\Large{$x$}}
\uput[d](0,8){\Large{$y$}}
\end{pspicture}
}
\end{center}
\end{figure}  
\begin{enumerate}[noitemsep, label=\textbf{\alph*}. ] 
\item 
What are the lengths of the opposite sides of $FGHI$?
item Are the opposite sides of $FGHI$ parallel?
\item  Do the diagonals of $FGHI$ bisect each other?
\item  Can you state what type of quadrilateral $FGHI$ is? Give reasons for your answer.
\end{enumerate}
\item 
A quadrialteral $ABCD$ with vertices $A(3;2)$, $B(1;7)$, $C(4;5)$ and $D(1;3)$ is given.
\begin{enumerate}[noitemsep, label=\textbf{\alph*}. ] 
\item  Draw the quadrilateral.
\item  Find the lengths of the sides of the quadrilateral.
\end{enumerate}
\item $ABCD$ is a quadrilateral with verticies $A(0;3)$, $B(4;3)$, $C(5;-1)$ and $D(-1;-1)$.
\begin{enumerate}[noitemsep, label=\textbf{\alph*}. ] 
\item Show that:
\begin{enumerate}[noitemsep, label=\textbf{\roman*}. ] 
\item $AD = BC$
\item $AB \parallel DC$
\end{enumerate}
\item What name would you give to $ABCD$?
\item Show that the diagonals $AC$ and $BD$ do not bisect each other.
\end{enumerate}
\item $P$, $Q$, $R$ and $S$ are the points $(-2;0)$, $(2;3)$, $(5;3)$, $(-3;-3)$ respectively.
\begin{enumerate}[noitemsep, label=\textbf{\alph*}. ] 
\item Show that:
\begin{enumerate}[noitemsep, label=\textbf{\roman*}. ] 
\item $SR = 2PQ$
\item $SR \parallel PQ$
\end{enumerate}
\item Calculate:
\begin{enumerate}[noitemsep, label=\textbf{\roman*}. ] 
\item $PS$
\item $QR$
\end{enumerate}
\item What kind of a quadrilateral is $PQRS$? Give reasons for your answers.
\end{enumerate}
\item $EFGH$ is a parallelogram with verticies $E(-1;2)$, $F(-2;-1)$ and $G(2;0)$. Find the coordinates of H by using the fact that the diagonals of a parallelogram bisect each other.\newline
\item  
$PQRS$ is a quadrilateral with points $P(0;-3)$ ; $Q(-2;5)$ ; $R(3;2)$ and $S(3;-2)$  in the Cartesian plane.
\begin{enumerate}[noitemsep, label=\textbf{\alph*}. ] 
\item Find the length of $QR$.
\item Find the gradient of $PS$.
\item Find the mid-point of $PR$.
\item Is $PQRS$ a parallelogram?  Give reasons for your answer.
 \end{enumerate}
\item $A(-2;3)$ and $B(2;6)$ are points in the Cartesian plane. $C(a;b)$ is the mid-point of $AB$. Find the values of $a$ and $b$.
\item 
Consider: Triangle $ABC$ with vertices $A(1; 3)$ , $B(4;1)$ and $C (6; 4)$:
\\begin{enumerate}[noitemsep, label=\textbf{\alph*}. ] 
\item Sketch triangle $ABC$ on the Cartesian plane. 
\item Show that $ABC$ is an isoceles triangle.
\item Determine the coordinates of $M$, the mid-point of $AC$.
\item Determine the gradient of $AB$.
\item Show that the following points are collinear: $A$, $B$ and $D(7;-1)$
\end{enumerate}
\item In the diagram, $A$ is the point $(-6;1)$ and $B$ is the point $(0;3)$
\setcounter{subfigure}{0}
\begin{figure}[H] % horizontal\label{m39107*id63458}
\begin{center}
\scalebox{.8}{
\begin{pspicture}(-5,-5)(5.5,5.5)
% \psaxes{<->}(0,0)(5,5)
\psgrid[gridcolor=lightgray,linecolor=lightgray,subgriddiv=1](0,0)(-10,-1)(1,7)
\psaxes[linewidth=1pt,labels=none,ticks=none]{->}(0,0)(-10,-1)(1,7)
\psline[linewidth=.05cm](-6,1)(0,3)
% \psline[linewidth=.05cm](0,0)(3,3.5)
% \psline[linewidth=.05cm](0,0)(3,2)
% \psline[linewidth=.05cm](0,0)(3,0.5)
% \uput[ur](.9,3.5){\Large{$T$}}
\uput[d](-6,1){\Large{$A (-6,1)$}}
\uput[d](-5.9,1.2){\qdisk(0,0){3pt}}
\uput[ul](-.2,3){\Large{$B (0,3)$}}
\uput[d](0,3.2){\qdisk(0,0){3pt}}
\uput[l](1.8,0){\Large{$x$}}
\uput[d](0,8){\Large{$y$}}
\end{pspicture}
}
\end{center}
\end{figure} 
\\begin{enumerate}[noitemsep, label=\textbf{\alph*}. ] 
\item Find the equation of line $AB$ 
\item Calculate the length of $AB$
\item  $A'$ is the image of $A$ and $B'$ is the image of $B$. Both these images are obtain by applying the transformation: $(x;y)\to(x - 4;y - 1)$. Give the coordinates of both $A'$ and $B'$
\item Find the equation of $A'B'$
\item Calculate the length of $A'B'$
\item Can you state with certainty that $AA'B'B$ is a parallelogram? Justify your answer.
\end{enumerate}
\end{enumerate}
\
\par \raisebox{-5 pt}{\includegraphics[width=0.5cm]{col11306.imgs/summary_www.png}} Find the answers with the shortcodes:
\par \begin{tabular}[h]{cccccc}
(1.) lgv  &  (2.) liZ  &  (3.) liB  &  (4.) lac  &  (5.) lax  &  (6.) laa  &  (7.) laY  &  (8.) lag  &  (9.) la4  &  (10.) l4o  & \end{tabular}
\end{eocexercises}